\documentclass[11pt,a4paper]{article}
\usepackage[utf8]{inputenc}
\usepackage[T1]{fontenc}
\usepackage[brazil]{babel}
\usepackage[margin=2.5cm]{geometry}
\usepackage{mathptmx}
\usepackage{enumitem}
\usepackage{booktabs}
\usepackage{tabularx}
\usepackage{array}
\usepackage{hyperref}
\hypersetup{colorlinks=false, pdfborder={0 0 0}}
\usepackage{fancyhdr}

\pagestyle{fancy}
\fancyhf{}
\fancyhead[L]{\small TrackFleet --- Plano de Comunicações}
\fancyhead[R]{\small\thepage}
\renewcommand{\headrulewidth}{0.4pt}

\newcolumntype{L}{>{\raggedright\arraybackslash}X}

\begin{document}

\begin{center}
{\Large\bfseries DOCUMENTO DE PARTES INTERESSADAS DO PROJETO}\\[2pt]
{\large\bfseries TrackFleet --- Gestão de Rastreador de Automóveis}\\[4pt]
Disciplina de Gerenciamento de Projetos $\cdot$ Framework PMBOK\textregistered\ 8\textsuperscript{a} Edição $\cdot$ Domínio: Partes Interessadas\\[2pt]
Integrantes: Enzo Allebrand e Kerlon Ribeiro (dois GPs) $\cdot$ 4\textsuperscript{o} semestre $\cdot$ Data: 27/08
\end{center}

\vspace{2mm}
\hrule
\vspace{4mm}

\noindent\textbf{\large Plano de Comunicações}

\vspace{1mm}
\noindent Documento operacional que define o fluxo de informações do projeto: quem precisa de qual informação, em que formato, por qual canal e com qual frequência, garantindo que as comunicações sejam personalizadas para cada grupo.

\section{Matriz de Comunicações do Projeto}

\begin{table}[h]
\centering
\small
\begin{tabularx}{\textwidth}{@{} L l l l l @{}}
\toprule
\textbf{O quê (Informação)} & \textbf{Para quem} & \textbf{Formato e Canal} & \textbf{Frequência} & \textbf{Responsável} \\
\midrule
\textbf{Relatório Executivo de Status} (Avanço, problemas, custos) & Diretoria (Patrocinador) & Documento PDF (1 pág.) enviado via E-mail & Quinzenal & GPs \\
\addlinespace
\textbf{Demonstração e Validação} (Entregas e protótipos) & Gestores de Frota & Reunião online (Teams/Meet) com atas pós-reunião & Mensal (ao fim da iteração) & GPs \\
\addlinespace
\textbf{Avisos de Privacidade e Uso} & Motoristas & Informativo impresso e termos de uso no app & Única / Em atualizações & Transportadora \\
\addlinespace
\textbf{Sincronização Técnica} (Bloqueios, decisões de código) & Equipe (Enzo e Kerlon) & Reunião rápida via Discord e registro no GitHub & Diária & Ambos os GPs \\
\bottomrule
\end{tabularx}
\end{table}

\section{Monitoramento das Comunicações}

\begin{itemize}[itemsep=4pt]
    \item \textbf{Coleta de Feedback:} A cada final de mês, será solicitado feedback informal ao Patrocinador para avaliar se o relatório executivo está no nível de detalhe correto (sem excesso de termos técnicos).
    \item \textbf{Ajuste de Rota:} Caso a equipe identifique falhas de alinhamento técnico entre os GPs, a frequência de sincronização ou a ferramenta de registro de código (GitHub) será revisada e adaptada conforme a necessidade.
\end{itemize}

\end{document}
